\section{Introduction}

The first wall of a fusion reactor is a two-material system. A plasma-facing
material meets the plasma, withstanding extreme heat flux, particle bombardment,
and neutron irradiation, while a structural alloy behind it carries the mechanical
load. Tungsten is the leading plasma-facing material owing to its high melting
point, low sputtering yield, and good thermal conductivity, but it is brittle and
unsuited to a load-bearing role~\cite{tungstenPFMreview2023}. Reduced-activation
ferritic-martensitic (RAFM) steels such as EUROFER97 serve as the structural
material, with favorable activation but limited high-temperature
capability~\cite{lindau2005eurofer,RAFMreview2017}. Both must also resist
irradiation damage and remain acceptably activated through operation and
decommissioning~\cite{zinkle2000operating,bloom2007critical}. As no single material
satisfies the surface and structural requirements simultaneously, the first wall
must be built from a plasma-facing material and a structural material joined into
one component.

Joining these materials is a central difficulty. Tungsten and steel differ
markedly in thermal expansion, with coefficients of roughly \(4.4\) and
\(12.7 \times 10^{-6}\)~K\(^{-1}\), so an abrupt joint concentrates large thermal
stresses, and direct tungsten--iron contact forms brittle intermetallics,
principally Fe\(_2\)W and Fe\(_7\)W\(_6\), that embrittle the
joint~\cite{fgmWsteelCGS2022,karim2024wire}. Functionally graded materials (FGMs),
in which composition or structure varies with position, offer a potential solution.
A compositionally graded alloy (CGA), a subset of FGMs in which the alloy
composition varies gradually from one terminal alloy to another, can spread the
expansion mismatch over a graded region rather than a single
boundary~\cite{fgmWEurofer2015}. Such gradients are now practically realizable, as
additive manufacturing (AM) processes such as directed energy deposition can vary
the feed composition layer by layer~\cite{karim2024wire}.

Despite these advances in manufacturing capability, grading alone does not avoid
the intermetallics, since iron and tungsten still meet across the gradient. One
remedy is to keep the transition a single body-centered cubic (BCC) solid solution
at every layer, so that the gradient never enters the composition ranges where the
brittle phases are stable~\cite{fgmWsteelInterlayerBCC2023}. This choice changes
the structural terminal itself. Rather than grading toward a steel, the gradient can
terminate in a vanadium-rich structural alloy that is reduced in activation and
reachable from tungsten through a continuous BCC path, so that iron and tungsten
never meet along the transition. Existing realizations, however, rely on a few
interlayers selected by hand, one composition at a
time~\cite{fgmWsteelInterlayerBCC2023}. Identifying such a path by inspection is a
fundamentally limited approach, since the constraints interact across a
composition space far too large to visualize.

The constraints are also numerous. The plasma-facing side requires high-temperature
capability and the structural side requires ductility, while reduced activation, a
stable single-phase structure, adequate thermal conductivity, and sufficient
strength must be met by both terminals and every intermediate layer alike.
Multicomponent refractory alloys offer wide freedom to balance these
demands~\cite{HEAnuclear2021,Senkov2010}, and high-throughput screening that
combines CALculation of PHAse Diagrams (CALPHAD) thermodynamics with physics-based
strength and ductility models has proven effective in reducing such spaces to
actionable candidate sets~\cite{vela2023data,HTcalphad2022}. Among the criteria,
neutron-induced activation is especially consequential, since transmutation governs
decay heat, dose rates, maintenance, and waste burden~\cite{RAFMreview2017}. The
FLARE model makes this criterion tractable by expressing the activation response of
any alloy as a composition-weighted sum of precomputed elemental
responses~\cite{Marquez2026FLARE}. Because these responses are linear in
composition, the compositions satisfying all activation limits form a convex
polytope that can be obtained directly by linear programming, without evaluating
the candidates one by one.

Computational path planning offers a systematic alternative to hand-selected
interlayers. Kirk et al.\ adapted motion-planning algorithms from robotics to plan
gradient paths through high-dimensional composition spaces while avoiding regions
of undesirable phases~\cite{kirk2018gradient}, an approach subsequently
demonstrated in additively manufactured gradients~\cite{eliseeva2019fgm}. The
framework was later extended to paths with monotonic property profiles, using a
lack-of-monotonicity cost function, on the argument that a monotonic gradient can
be transformed into nearly any spatial profile on the part by controlling the
deposition rate~\cite{kirk2021monotonic}. This cost function has since been used to
design and print a gradient that is monotonic in stacking fault
energy~\cite{hanagan2025sfe}, and monotonic property profiles have also been
imposed in graded deposition, where a liquidus that falls along the build prevents
the remelting of previously deposited layers~\cite{glerum2026ss316tzm}. More
recently, the composition space has been represented as a graph of feasible
compositions, enabling the design of multi-terminal graded
structures~\cite{multiterminalCGA2024} and the identification of connected feasible
regions for CGA design~\cite{obrist2026graph}. These methods, however, have treated
monotonicity in a single property at a time. The first wall requires two at once:
a coefficient of thermal expansion (CTE) that never decreases, so that the
expansion mismatch at every interface acts in the same direction, and a solidus
that never increases, so that no layer melts at a lower temperature than a layer
behind it, which faces a lower temperature in service. Both must hold along a path
whose every composition also satisfies activation, phase-stability, conductivity,
and strength constraints.

In this work, a CGA for the fusion first wall is designed in the eleven-element
refractory system \ch{Cr-Fe-Hf-Mo-Nb-Re-Ta-Ti-V-W-Zr}. Activation screening by
linear programming, referenced to tungsten and EUROFER97, defines the feasible
composition space. CALPHAD calculations reduce this space to a connected,
single-phase BCC design region, whose compositions are classified as plasma-facing
or structural materials by solidus, Pugh ratio, and EUROFER97-level yield strength.
The gradient path between the two terminal alloys is then planned on the
composition graph with the cost function of Kirk et al.~\cite{kirk2021monotonic},
with monotonicity in both CTE and solidus imposed as hard constraints. Under these
constraints the cost reduces to path length, and the optimal path is found exactly
by Dijkstra's algorithm, with no penalty weight to tune. The approach improves upon
previous work by enabling monotonicity in more than one property at once while
every layer remains reduced in activation, single-phase, and at least as strong as
EUROFER97.

The remainder of this paper is organized as follows. First, the activation
screening, property calculations, and gradient path design problem are described in
Section~\ref{sec:methods}. The resulting feasible space, design region, terminal
alloys, and gradient path are then presented and discussed in
Section~\ref{sec:results}. Finally, conclusions and future work are given in
Section~\ref{sec:conclusion}. The overall workflow is summarized in
Fig.~\ref{fig:overview}.

\begin{figure*}[!t]
  \centering
  \includegraphics[width=0.9\textwidth]{figures/FGM-SCHEMATIC.png}
 \caption{Overview of the design workflow: activation screening by linear
programming, property filtering and classification of the feasible composition
space, selection of the plasma-facing and structural terminal alloys, and gradient
path design with monotonic CTE and solidus.}
  \label{fig:overview}
\end{figure*}
